\section*{Capítulo \ref{ch:g4:decimals} --- Décimas y centésimas}

\begin{solution}{exo:g4:decimals:1}
$5$ columnas: $\frac{5}{10}$ (es decir, $\frac{50}{100}$). $2$ columnas
y $3$ casillas: $\frac{23}{100}$. $45$ casillas: $\frac{45}{100}$.
\end{solution}

\begin{solution}{exo:g4:decimals:2}
$0.9$; \quad $0.27$; \quad $0.03$; \quad $5.6$; \quad
$0.60 = 0.6$.
\end{solution}

\begin{solution}{exo:g4:decimals:3}
$0.3 = \frac{3}{10}$; \quad $0.51 = \frac{51}{100}$; \quad
$0.07 = \frac{7}{100}$; \quad $1.9 = 1 + \frac{9}{10} =
\frac{19}{10}$.
\end{solution}

\begin{solution}{exo:g4:decimals:4}
El $5$ cuenta las décimas y el $8$, las centésimas:
\[
6.58 = 6 + \frac{5}{10} + \frac{8}{100}.
\]
\end{solution}

\begin{solution}{exo:g4:decimals:5}
$4.2$: en la segunda marca después del $4$. $4.75$: entre las marcas
$4.7$ y $4.8$, a \omterm{def:g3:division:half}{mitad} de camino. $4.9$: una marca antes del $5$.
\end{solution}

\begin{solution}{exo:g4:decimals:6}
$0.6 = 0.60$; \quad $2.8 > 2.75$ (compara $2.80$); \quad
$0.09 < 0.1$ ($9$ centésimas frente a $10$); \quad
$5.3 > 5.13$.
\end{solution}

\begin{solution}{exo:g4:decimals:7}
$0.95 < 1.05 < 1.45 < 1.5$.
\end{solution}

\begin{solution}{exo:g4:decimals:8}
$3.05$; \quad $12.50$; \quad $0.99$. Orden:
$0.99 < 3.05 < 12.50$.
\end{solution}

\begin{solution}{exo:g4:decimals:9}
$3.1$, $3.2$, $3.3$, $3.4$, $3.5$, $3.6$, $3.7$, $3.8$, $3.9$ ---
nueve números.
\end{solution}

\begin{solution}{exo:g4:decimals:10}
En el cuadrado de cien, $0.25$ sombrea $25$ casillas, mientras que $0.5 = 0.50$
sombrea $50$: el mayor es $0.5$. Tomás comparó $25$ y
$5$ como números enteros; tenía que comparar $25$ y $50$ centésimas.
\end{solution}

\begin{solution}{exo:g4:decimals:11}
Orden: $3.58 < 3.6 < 3.65$. Los saltos $3.6$ y $3.65$ se diferencian en
$0.05 = \frac{5}{100}$ de metro (cinco centímetros).
\end{solution}
